% ============================================================
%  第七章  基于结构的反事实逻辑（第 7.1—7.2 节）
% ============================================================
\chapter{基于结构的反事实逻辑}

\begin{flushright}
\kaishu\small
主说：“我若在所多玛城里见有五十个义人，我就为他们的缘故，饶恕那地方的众人。”\\[2pt]
——《创世记》18:26
\end{flushright}

\section*{序言}
\addcontentsline{toc}{section}{序言}

本章对基于结构的反事实（structure-based counterfactual）作形式分析。这一概念在第一、三章已简要引入，本书其后的讨论将一直围绕它展开。通过这一分析，我们为此前各章引入的其他概念——因果模型、行动、因果效应、因果相关性（causal relevance）、误差项与外生性——取得更精确的数学定义。

在把因果模型与反事实的概念表述为形式化的数学对象之后，我们将用例子演示：反事实问题如何从确定性的与概率性的因果模型中都能得到回答（\S 7.1）。在 \S 7.2.1 中，我们将论证政策分析乃是一种反事实推理，并以计量经济学中的一个简单例子演示这一论点。这将为 \S 7.2.2 的讨论作好铺垫——那里我们用政策预测来阐明反事实的经验内容。\S 7.2.3 讨论反事实在因果解释的诠释与生成中所起的作用。\S 7.2 以两项讨论收尾：因果关系如何从行动与机制中浮现（\S 7.2.4），以及因果方向性如何能从一组对称方程中被诱导出来（\S 7.2.5）。

在 \S 7.3 中，我们给出反事实关系与因果相关关系的一个公理化刻画，这些关系正是从结构模型语义中浮现出来的。\S 7.3.1 将识别出一组性质（即公理），使我们能从假设推出新的反事实关系；\S 7.3.2 演示这些公理在因果效应的代数推导中的用法；\S 7.3.3 引入因果相关关系的公理，并借助它们与图的公理之间的相似性，说明如何用图来验证相关性关系。

\S 7.3 所发展的公理化刻画，使我们能把结构模型与处理因果性和反事实的其他进路相比较，其中最值得注意的是基于 Lewis 最近世界语义的进路（\S 7.4.1—7.4.4）。结构进路与 Neyman–Rubin 潜在结局框架的形式等价性，在 \S 7.4.4 中讨论。最后，我们重温外生性这一主题，并在 \S 7.4.5 中用外生变量与工具变量的反事实定义，扩展 \S 5.4.3 的讨论。

本章的最后一部分（\S 7.5）把因果性的结构说明与基于概率关系的说明相比较。我们详述自己何以更倾向结构说明，并着重指出概率说明目前面临的困难。

\section{结构模型语义}

科学家如何从在完全不同的条件下进行的其他实验的结果，来预测某一次实验的结果？这类预测要求我们设想：在各种假设性的改变之下，世界会是什么样子；因而它诉诸反事实推断。反事实推断虽属科学思维的基本成分，却难以在逻辑、代数方程或概率的标准语言中得到形式化。反事实推断的形式化需要这样一种语言：在其中，世界中不变的关系得以与表示我们对世界之信念的暂时关系相区分；而标准代数——包括方程代数、布尔代数与概率演算——都不支持这种区分。结构模型提供了这种区分。本节给出反事实的结构模型语义，其定义见于 Balke and Pearl (1994a,b)、Galles and Pearl (1997, 1998) 与 Halpern (1998)，\footnote{类似的模型（称为“神经元图”，neuron diagram；Lewis 1986, p.~200; Hall 2004）被哲学家非正式地用来例示因果过程的链条。}它与 Rubin (1974) 的实验视角形成鲜明对照。Simon and Rescher (1966) 与 Ortiz (1999) 也提出过相关的进路。

我们从因果模型的一个确定性定义开始。如前面各章所论，该定义由所关心变量之间的函数关系构成，每个关系代表一个自主的机制。因果关系与反事实关系，在此模型中被定义为对这些机制作局部修改时所引起的响应。概率关系则通过对背景条件赋予概率而自然浮现。在用例子演示这一模型如何便利于确定性与概率性语境下的反事实计算（\S 7.1.2）之后，我们给出一套借助因果图计算反事实表达式之概率的一般方法（\S 7.1.3）。

\subsection{定义：因果模型、行动与反事实}

“模型”一词在日常用法中指对现实的理想化表示，它凸显某些方面而忽略另一些方面。然而在逻辑系统中，模型是一个数学对象，它为给定语言中的语句赋予真值，其中每条语句表示现实的某个方面。例如，真值表是命题逻辑中的模型：它为任一布尔表达式赋予真值，而该表达式可以表示所论领域中的一个事件或一组条件。又如，联合概率函数是概率逻辑中的模型：它为形如 $P(A\mid B)<p$ 的任何语句赋予真值，其中 $A$ 与 $B$ 是表示事件的布尔表达式。

因果模型自然应当编码那些涉及因果关系的语句的真值；这类语句包括行动语句（如“若我们做 $B$，则 $A$ 将为真”）、反事实（如“若不是因为 $B$，$A$ 本会有所不同”）以及素朴的因果话语（如“$A$ 可能导致 $B$”或“$B$ 之所以发生是因为 $A$”）。此类语句无法在标准命题逻辑或概率演算中得到解释，因为它们所处理的是外部世界中发生的变化，而不是我们对一个静态世界的信念的变化。因果模型通过对这些变化中被改动的机制作出显式表示，来编码并区分关于外部变化的信息。

\begin{definition}[因果模型]
一个\emph{因果模型}（causal model）是一个三元组
\[
M=\langle U,V,F\rangle,
\]
其中：
\begin{enumerate}[label=(\roman*),leftmargin=2.6em]
\item $U$ 是一组\emph{背景变量}（background variable，亦称外生变量），\footnote{我们将尽量避免用“外生”一词指称背景条件，因为该词已带有更精细的技术含义（见 \S 5.4.3 与 \S 7.4）。计量经济学文献中常用“预定”（predetermined）一词。}它们由模型之外的因素决定；
\item $V$ 是变量集合 $\{V_1,V_2,\dots,V_n\}$，称为\emph{内生变量}，它们由模型中的变量决定——即由 $U\cup V$ 中的变量决定；且
\item $F$ 是函数集合 $\{f_1,f_2,\dots,f_n\}$，其中每个 $f_i$ 是从（$U_i\cup\PA_i$ 各自的定义域）到 $V_i$ 的映射，这里 $U_i\subseteq U$、$\PA_i\subseteq V\backslash V_i$，且整个集合 $F$ 构成从 $U$ 到 $V$ 的映射。换言之，$F$ 中的每个 $f_i$ 为 $V_i$ 赋一个依赖于 $V\cup U$ 中某一精选变量组（之取值）的值：
\[
v_i=f_i(\pa_i,u_i),\qquad i=1,\dots,n,
\]
而整个集合 $F$ 有唯一解 $V(u)$。\footnote{$\PA_i$（意为“父节点”，parents）的选取并非任意，而是表达了建模者的如下理解：在决定 $V_i$ 的取值之前，“自然”必须查询哪些变量。}\footnote{唯一性在递归（即无环）系统中得到保证。Halpern (1998) 允许非递归系统有多重解。}
\end{enumerate}
\end{definition}

每个因果模型 $M$ 都可联系一个有向图 $G(M)$：其中每个节点对应一个变量，有向边由 $\PA_i$ 与 $U_i$ 的成员指向 $V_i$。我们称这样的图为与 $M$ 相联系的\emph{因果图}。此图仅仅标明哪些内生变量与背景变量直接影响各个 $V_i$，并不指定 $f_i$ 的函数形式。把父节点集 $\PA_i$ 限于 $V$ 中变量这一惯例，源于背景变量往往不可观测。不过一般说来，我们可以把父节点集扩展到包含 $U$ 中的可观测变量。

\begin{definition}[子模型]
设 $M$ 为因果模型，$X$ 为 $V$ 中的一组变量，$x$ 为 $X$ 的一个特定取值。$M$ 的\emph{子模型}（submodel）$M_x$ 是如下因果模型
\[
M_x=\langle U,V,F_x\rangle,
\]
其中
\begin{equation}
F_x=\{f_i: V_i\notin X\}\cup\{X=x\}. \tag{7.1}
\end{equation}
\end{definition}

用文字说，$F_x$ 是这样构成的：从 $F$ 中删去与集合 $X$ 的成员相对应的全部函数 $f_i$，并以常数函数集 $X=x$ 取而代之。

子模型适于表示局部行动与假设性改变（包括反事实前件所蕴含的那些改变）的效果。若把 $F$ 中每个函数 $f_i$ 解释为一个独立的物理机制，并把行动 $\dox{X=x}$ 定义为“为使 $X=x$ 在任何 $u$ 下成立而对 $M$ 所需作的最小改动”，则 $M_x$ 就表示这一最小改动所得的模型，因为它与 $M$ 的差别仅在于那些直接决定 $X$ 中变量的机制。从 $M$ 到 $M_x$ 的变换修改了 $F$ 的代数内容；这正是 Galles and Pearl (1998) 使用“可修改结构方程”（modifiable structural equations）这一名称的缘由。\footnote{结构性的修改可上溯至 Haavelmo (1943)、Marschak (1950) 与 Simon (1953)。把干预显式翻译为“抹去”模型中的方程，最早由 Strotz and Wold (1960) 提出，其后为 Fisher (1970)、Sobel (1990)、Spirtes et al. (1993) 与 Pearl (1995a) 所用。Fine (1985) 引入过类似的子模型概念，但并非专为表示行动与反事实。}

\begin{definition}[行动的效果]
设 $M$ 为因果模型，$X$ 为 $V$ 中的一组变量，$x$ 为 $X$ 的一个特定取值。行动 $\dox{X=x}$ 对 $M$ 的\emph{效果}由子模型 $M_x$ 给出。
\end{definition}

\begin{definition}[潜在响应]
设 $X$ 与 $Y$ 是 $V$ 中两个变量子集。$Y$ 对行动 $\dox{X=x}$ 的\emph{潜在响应}（potential response），记作 $Y_x(u)$，是方程组 $F_x$ 关于 $Y$ 的解，\footnote{若 $Y$ 是变量集 $Y=(Y_1,Y_2,\dots)$，则 $Y_x(u)$ 表示函数向量 $(Y_{1x}(u),Y_{2x}(u),\dots)$。}即
\[
Y_x(u)=Y_{M_x}(u).
\]
\end{definition}

我们将把注意力限于 $\dox{X=x}$ 形式的行动。形如“若 $Z=z$ 则 $\dox{X=x}$”的条件行动，可以用“以 $Z$ 的函数而非常数替换方程”的方式来形式化（\S 4.2）。形如“$\dox{X=x}$ 或 $Z=z$”的析取式行动则不加考虑，因为它们会使反事实的概率处理复杂化。

\begin{definition}[反事实]
设 $X$ 与 $Y$ 是 $V$ 中两个变量子集。反事实语句“（在情形 $u$ 下）若 $X$ 曾为 $x$，则 $Y$ 会为 $y$”被解释为等式
\[
Y_x(u)=y,
\]
其中 $Y_x(u)$ 是 $Y$ 对 $X=x$ 的潜在响应。
\end{definition}

定义 7.1.5 于是用对方程的假设性修改，来解释反事实短语“若 $X$ 曾为 $x$”：它模拟一个外部行动（或自发的变化），该行动修改了实际的历史进程，并以对机制的最小改动强制条件“$X=x$”成立。这是反事实语义中关键的一步（Balke and Pearl 1994b），因为它允许 $x$ 不同于 $X(u)$ 的当前取值而不造成逻辑矛盾；它同时也抑制了从反事实前件 $X=x$ 出发的溯因推断（或回溯）。\footnote{Simon and Rescher (1966, p.~339) 在他们关于反事实的说明中没有包含这一步，并指出由前件触发的反向推断可能导致含混的解释。}在第三章（\S 3.6.3）中，我们用记号 $Y(x,u)$ 表示基于个体的条件句“在个体 $u$ 中、若 $X$ 曾为 $x$，$Y$ 将会取得的值”，即 Neyman–Rubin 潜在结局模型中的用法。本书其余部分将用记号 $Y_x(u)$ 表示与定义 7.1.5 的结构模型解释特别相联系的反事实（与式 \textup{(3.51)} 平行）；$Y(x,u)$ 则保留给一般性的虚拟条件句，不承诺任何特定的语义。

定义 7.1.5 赋予原子机制 $\{f_i\}$ 本身以干预—反事实的解释，因为 $v_i=f_i(\pa_i,u_i)$ 就是 $V_i$ 在子模型 $M_{V\backslash v_i}$ 中的取值。换言之，$f_i(\pa_i,u_i)$ 表示当我们对 $V$ 中所有其他变量都施行干预时 $V_i$ 的潜在响应。

这一表述可自然地推广到概率系统，如下。

\begin{definition}[概率因果模型]
一个\emph{概率因果模型}（probabilistic causal model）是一个二元组
\[
\langle M,P(u)\rangle,
\]
其中 $M$ 是因果模型，$P(u)$ 是定义在 $U$ 的定义域上的概率函数。
\end{definition}

函数 $P(u)$，连同“每个内生变量都是 $U$ 的函数”这一事实，定义了内生变量上的一个概率分布。也就是说，对任一变量集 $Y\subseteq V$，有
\begin{equation}
P(y)\triangleq P(Y=y)=\sum_{\{u\mid Y(u)=y\}}P(u). \tag{7.2}
\end{equation}
反事实语句的概率以同样的方式、通过子模型 $M_x$ 所诱导的函数 $Y_x(u)$ 来定义：
\begin{equation}
P(Y_x=y)=\sum_{\{u\mid Y_x(u)=y\}}P(u). \tag{7.3}
\end{equation}
同样，因果模型在反事实语句上定义了一个联合分布。也就是说，$P(Y_x=y,Z_w=z)$ 对任意（不必互不相交的）变量集 $Y,X,Z,W$ 都有定义。特别地，$P(Y_x=y,X=x')$ 与 $P(Y_x=y,Y_{x'}=y')$ 对 $x\ne x'$ 都有良好定义，且分别由
\begin{equation}
P(Y_x=y,X=x')=\sum_{\{u\mid Y_x(u)=y\ \&\ X(u)=x'\}}P(u) \tag{7.4}
\end{equation}
与
\begin{equation}
P(Y_x=y,Y_{x'}=y')=\sum_{\{u\mid Y_x(u)=y\ \&\ Y_{x'}(u)=y'\}}P(u) \tag{7.5}
\end{equation}
给出。若 $x$ 与 $x'$ 不相容，则 $Y_x$ 与 $Y_{x'}$ 无法同时被测量，为联合语句“若 $X=x$ 则 $Y$ 会为 $y$，且若 $X=x'$ 则 $Y$ 会为 $y'$”赋予概率，似乎没有意义。这类顾虑正是近来反对把反事实当作联合分布的随机变量来处理的一个来源（Dawid 2000）。然而，$Y_x$ 与 $Y_{x'}$ 是作为 $Y$ 在两个不同子模型中的解来定义的，而这两个子模型都由 $U$ 上的标准概率空间支配；这一来，上述顾虑便被消解了——那个自相矛盾的联合语句，被解释为 $U$ 空间中的一个普通事件。

我们特别感兴趣的是以实际观测为条件的反事实概率。例如，事件 $X=x$ “是”事件 $Y=y$ 之原因的概率，可以解释为：在已知 $X=x$ 与 $Y=y$ 确实发生的条件下，若 $X$ 不曾为 $x$，则 $Y$ 不会等于 $y$ 的概率（关于因果概率的深入讨论见第九章）。这类概率在上述模型中都有良好定义；它们要求对形如 $P(Y_{x'}=y'\mid X=x,Y=y)$ 的表达式求值，其中 $x'$ 与 $y'$ 分别与 $x$ 与 $y$ 不相容。式 \textup{(7.4)} 使我们能如下计算该量：
\begin{equation}
\begin{aligned}
P(Y_{x'}=y'\mid X=x,Y=y) &= \frac{P(Y_{x'}=y',X=x,Y=y)}{P(X=x,Y=y)}\\
&= \sum_{u}P(Y_{x'}(u)=y')P(u\mid x,y).
\end{aligned} \tag{7.6}
\end{equation}
换言之，我们先把 $P(u)$ 更新为 $P(u\mid x,y)$，再用更新后的分布 $P(u\mid x,y)$ 计算等式 $Y_{x'}(u)=y'$ 的期望。这就为 \S 1.4 引入的三步程序提供了依据；我们现将其总结为一条定理。

\begin{theorem}
给定模型 $\langle M,P(u)\rangle$，反事实语句“若当时为 $A$，则 $B$”在给定证据 $e$ 下的条件概率 $P(B_A\mid e)$，可用以下三步求值。
\begin{enumerate}[label=\arabic*.,leftmargin=2.2em]
\item \textbf{溯因}（abduction）——用证据 $e$ 更新 $P(u)$，得到 $P(u\mid e)$。
\item \textbf{行动}（action）——用行动 $\dox{A}$（$A$ 为反事实的前件）修改 $M$，得到子模型 $M_A$。
\item \textbf{预测}（prediction）——用修改后的模型 $\langle M_A,P(u\mid e)\rangle$ 计算 $B$（反事实的后件）的概率。
\end{enumerate}
\end{theorem}

为结束本节，我们再引入两个对象，它们在后续讨论中将证明有用：世界\footnote{Adnan Darwiche 提请我注意这一对象的重要性。}与理论。

\begin{definition}[世界与理论]
一个\emph{因果世界}（causal world）$w$ 是一个二元组 $\langle M,u\rangle$，其中 $M$ 是因果模型，$u$ 是背景变量 $U$ 的一个特定取值。一个\emph{因果理论}（causal theory）是若干因果世界构成的集合。
\end{definition}

世界 $w$ 可视为 $P(u)=1$ 的退化概率模型。因果理论将用来刻画因果模型的部分设定，例如共享同一因果图的模型，或 $f_i$ 为系数待定的线性函数的模型。

\subsection{反事实的求值：确定性分析}

在 \S 1.4.1 中，我们用若干例子演示了结构模型中行动与反事实的解释。现在我们施用 \S 7.1.1 的定义，演示确定性的与概率性的反事实查询如何能用结构模型语义形式地得到回答。

\subsubsection*{例 1：行刑队}

考虑图 7.1 所示的双人行刑队，其中 $A,B,C,D,U$ 分别代表以下命题：
\begin{align*}
U &= \text{法庭下令处决；}\\
C &= \text{队长发出信号；}\\
A &= \text{枪手 $A$ 开枪；}\\
B &= \text{枪手 $B$ 开枪；}\\
D &= \text{囚犯死亡。}
\end{align*}
假定法庭的决定未知，两名枪手都精准、警觉且守法，囚犯也不至于因惊恐或其他外因而死亡。我们希望构造这个故事的某个形式表示，使下列语句可以被机械地求值。

\begin{enumerate}[label=\textbf{S\arabic*},leftmargin=3.4em,itemsep=0.55em,topsep=0.4em]
\item \emph{预测}——若枪手 $A$ 没有开枪，则囚犯还活着：
\[
\neg A\implies \neg D.
\]
\item \emph{溯因}——若囚犯还活着，则队长没有发出信号：
\[
\neg D\implies \neg C.
\]
\item \emph{转导}（transduction）——若枪手 $A$ 开了枪，则 $B$ 也开了枪：
\[
A\implies B.
\]
\item \emph{行动}——若队长没有发出信号，而枪手 $A$ 决定开枪，则囚犯会死，且 $B$ 不会开枪：
\[
\neg C\implies D_A\ \&\ \neg B_A.
\]
\item \emph{反事实}——若囚犯已死，则即使枪手 $A$ 没有开枪，囚犯也会死：
\[
D\implies D_{\neg A}.
\]
\end{enumerate}

\begin{figure}[htbp]
\centering
\begin{tikzpicture}[causal,scale=1.05,every node/.style={transform shape}]
  \node[vtx] (U) at (0,2.75) {$U$};
  \node[vtx] (C) at (0,1.55) {$C$};
  \node[vtx] (A) at (-1.55,0.35) {$A$};
  \node[vtx] (B) at (1.55,0.35) {$B$};
  \node[vtx] (D) at (0,-0.95) {$D$};
  \draw[ed] (U)--(C);
  \draw[ed] (C)--(A);
  \draw[ed] (C)--(B);
  \draw[ed] (A)--(D);
  \draw[ed] (B)--(D);
  \node[font=\small,anchor=west] at (0.45,2.75) {（法庭下令）};
  \node[font=\small,anchor=west] at (0.45,1.55) {（队长）};
  \node[font=\small,anchor=west] at (1.98,0.35) {（枪手）};
  \node[font=\small,anchor=west] at (0.45,-0.95) {（死亡）};
\end{tikzpicture}
\caption{双人行刑队一例中的因果关系。}
\label{fig:7.1}
\end{figure}

\subsubsection*{标准语句的求值}

为证明前三条语句，我们无须诉诸因果模型；这些语句只涉及标准的逻辑联结词，因而可以用标准的逻辑演绎来处理。这个故事可以用任何便利的逻辑理论（一组命题语句）来刻画，例如
\[
T_1:\quad U\iff C,\quad C\iff A,\quad C\iff B,\quad A\vee B\iff D
\]
或
\[
T_2:\quad U\iff C\iff A\iff B\iff D,
\]
其中每个理论都容许两个逻辑模型
\[
m_1:\ \{U,C,A,B,D\}\qquad\text{与}\qquad m_2:\ \{\neg U,\neg C,\neg A,\neg B,\neg D\}.
\]
用文字说，任何刻画我们这个故事的理论 $T$，都应蕴含“五个命题或者全真、或者全假”；模型 $m_1$ 与 $m_2$ 把这两种可能显式地呈现出来。S1—S3 的有效性很容易验证：或者从 $T$ 推出，或者指出每条语句的前件与后件同属一个模型。

在着手分析 S4 与 S5 之前，有两点值得说明。第一，$T_1$ 与 $T_2$ 中的双向蕴含是支撑溯因所必需的；若改用单向蕴含（如 $C\implies A$），我们就无法从 $A$ 推出 $C$。在标准逻辑中，这种对称性抹去了预测（沿时间向前推理）、溯因（从证据推理到解释）与转导（从证据推理到解释、再从解释推理到预测）这三项任务之间的一切区别。使用双向蕴含时，这三种推理方式只在“对条件句前件与后件所赋予的解释”上不同，而不在推理方法上不同。在非标准逻辑（如逻辑程序设计）中，蕴含符号规定了推理的方向，甚至连逆否都不被许可，此时必须诉诸元逻辑的推理机制才能执行溯因（Eshghi and Kowalski 1989）。

第二，使 S1—S3 在标准逻辑中易于处理的特点在于：它们全都处理\emph{认知}推断——即从信念推理到关于一个静态世界的信念。例如，语句 S2 可以被阐明为：若我们发现囚犯还活着，我们便有理由相信队长没有发出信号。逻辑中的实质蕴含符号（$\implies$）并不越出这一狭窄含义；下文即将把这一点与反事实蕴含相对照。

\subsubsection*{行动语句的求值}

语句 S4 诉诸一个蓄意的行动：“枪手 $A$ 决定开枪”。按我们关于行动的讨论（例如见第四章或定义 7.1.3），任何此类行动都必定违反故事初始理论中的某些前提或机制。要形式地确定在行动之下什么保持不变，我们必须把因果关系纳入理论；仅凭逻辑关系是不够的。与我们的故事相对应的因果模型如下。

\begin{center}
\begin{tabular}{r@{\hspace{1.0em}}l}
\multicolumn{2}{c}{\textbf{模型 $M$}}\\[4pt]
 & $(U)$\\
$C=U$ & $(C)$\\
$A=C$ & $(A)$\\
$B=C$ & $(B)$\\
$D=A\vee B$ & $(D)$
\end{tabular}
\end{center}

这里我们使用等号而非蕴含号，是为了 (i) 允许双向推理；(ii) 强调每个方程（与逻辑语句不同）都代表一个自主的机制（用数据库的语言说，一条“完整性约束”）——除非被特意违反，它保持不变。我们还在每个方程旁附上括号内的符号，以显式标出该方程中的因变量（位于左端），从而表示与图 7.1 中各箭头相联系的因果不对称性。

为求 S4 的值，我们依定义 7.1.3 构造子模型 $M_A$：其中方程 $A=C$ 被 $A$ 取代（模拟枪手 $A$ 不顾信号而决定开枪）。

\begin{center}
\begin{tabular}{r@{\hspace{1.0em}}l}
\multicolumn{2}{c}{\textbf{模型 $M_A$}}\\[4pt]
 & $(U)$\\
$C=U$ & $(C)$\\
$A$ & $(A)$\\
$B=C$ & $(B)$\\
$D=A\vee B$ & $(D)$\\
\hline
\multicolumn{2}{l}{事实：$\neg C$}\\
\hline
\multicolumn{2}{l}{结论：$A,\,D,\,\neg B,\,\neg U,\,\neg C$}
\end{tabular}
\end{center}

我们看到，给定 $\neg C$，便可容易地推出 $D$ 与 $\neg B$，从而确认 S4 的有效性。

值得注意的是，像 S4 这样“有问题”的语句——其前件违反了故事的一条基本前提（即两名枪手都守法）——可以在讲述该故事的同一个确定性框架中自然地得到处理。传统逻辑学家与概率论者往往把 S4 这类语句斥为矛盾，并坚持用概率的方式重新表述问题，以便容忍对定律 $A=C$ 的例外。\footnote{我揣测，这一问题正是 1960 年代概率因果性兴起的主要推动力之一（综述见 \S 7.5）。}这类重新表述是不必要的；结构进路使我们能在其自然的确定性栖息地中处理日常的因果陈述，而无须先把它们浸入非确定性的装饰之中。在这一框架中，一切定律都被理解为“可废止的”（defeasible）默认表述——可因蓄意的干预而失效。物理学的基本定律当然依旧不可更改，但它们对任一给定情景的适用性，却可因施动者的行动或外部干预而被修改。

\subsubsection*{反事实的求值}

现在我们准备求反事实语句 S5 的值。依定义 7.1.5，反事实 $D_{\neg A}$ 表示 $D$ 在子模型 $M_{\neg A}$ 中的取值。这一取值是含混的，因为它依赖于 $U$ 的取值，而 $U$ 在 $M_{\neg A}$ 中并未被指定。观测 $D$ 消除了这种含混：一旦发现囚犯已死，我们便可推断法庭已下达命令（$U$），从而，若枪手 $A$ 当时没有开枪，则枪手 $B$ 会开枪并杀死囚犯，这就确认了 $D_{\neg A}$。

形式地，我们可以用定理 7.1.7 的步骤推出 $D_{\neg A}$（尽管此处不涉及概率）。我们先在原始模型 $M$ 中加入事实 $D$ 并求 $U$ 的值；然后构造子模型 $M_{\neg A}$，并用第一步求得的 $U$ 值重新求 $D$ 在 $M_{\neg A}$ 中的真值。这些步骤具体如下。

\begin{center}
\begin{tabular}{r@{\hspace{1.0em}}l}
\multicolumn{2}{c}{\emph{第 1 步}}\\[2pt]
\multicolumn{2}{c}{\textbf{模型 $M$}}\\[4pt]
 & $(U)$\\
$C=U$ & $(C)$\\
$A=C$ & $(A)$\\
$B=C$ & $(B)$\\
$D=A\vee B$ & $(D)$\\
\hline
\multicolumn{2}{l}{事实：$D$}\\
\hline
\multicolumn{2}{l}{结论：$U,\,A,\,B,\,C,\,D$}\\[16pt]
\multicolumn{2}{c}{\emph{第 2 步}}\\[2pt]
\multicolumn{2}{c}{\textbf{模型 $M_{\neg A}$}}\\[4pt]
 & $(U)$\\
$C=U$ & $(C)$\\
$\neg A$ & $(A)$\\
$B=C$ & $(B)$\\
$D=A\vee B$ & $(D)$\\
\hline
\multicolumn{2}{l}{事实：$U$}\\
\hline
\multicolumn{2}{l}{结论：$U,\,\neg A,\,C,\,B,\,D$}
\end{tabular}
\end{center}

注意，从第 1 步带到第 2 步的只有背景变量 $U$ 的取值；其他一切命题都必须在模型被新修改之后重新求值。这反映了如下理解：背景因素 $U$ 既不受模型中变量的影响，也不受机制 $\{f_i\}$ 的影响；因此，反事实的后件（在此即 $D$）必须在与实际世界所盛行的相同的背景条件下求值。事实上，背景变量正是信息从实际世界传向假设世界的主要载体；在把前者转变为后者的动态过程中，它们充当“不变性的守护者”（或曰持久性）。（这是 David Heckerman 的观察，私人通信。）

还要注意，这一求反事实的两步程序可以合并为一步。若我们用星号区分修改后变量与修改前变量，便可把 $M$ 与 $M_x$ 合并为一个逻辑理论，并通过该合并理论中的纯逻辑演绎来证明 S5 的有效性。为作说明，我们把 S5 写成 $D\implies D^*_{\neg A^*}$（读作：若 $D$ 在实际世界中为真，则在由修改 $\neg A^*$ 所创造的假设世界中 $D$ 也会为真），并在合并理论中如下证明 $D^*$ 的有效性。

\begin{center}
\begin{tabular}{r@{\hspace{1.0em}}l@{\hspace{2.2em}}r@{\hspace{1.0em}}l}
\multicolumn{4}{c}{\textbf{合并理论}}\\[4pt]
$C^*=U$ & & $C=U$ & $(C)$\\
$\neg A^*$ & & $A=C$ & $(A)$\\
$B^*=C^*$ & & $B=C$ & $(B)$\\
$D^*=A^*\vee B^*$ & & $D=A\vee B$ & $(D)$\\
\hline
\multicolumn{4}{l}{事实：$D$}\\
\hline
\multicolumn{4}{l}{结论：$U,\,A,\,B,\,C,\,D,\,\neg A^*,\,C^*,\,B^*,\,D^*$}
\end{tabular}
\end{center}

注意 $U$ 不必加星号，这反映了背景条件保持不变的假定。

此处值得反思 S4 与 S5 之间的差别。二者在句法上似乎完全相同，都涉及一个事实蕴含一个反事实，然而我们把 S4 称为“行动”语句，把 S5 称为“反事实”语句。差别在于给定的事实与反事实的前件（即“行动”部分）之间的关系。在 S4 中，给定的事实（$\neg C$）不受前件（$A$）的影响；在 S5 中，给定的事实（$D$）则可能受前件（$\neg A$）的影响。这两种情形的差别是根本性的，从它们的求值方法即可看出。求 S4 的值时，我们事先知道 $C$ 不会受模型修改 $\dox{A}$ 的影响；因此可以把 $C$ 直接加入修改后的模型 $M_A$。求 S5 的值时则相反，我们考虑的是从 $D$ 到 $\neg D$ 的一种可能反转，而这一反转应归因于修改 $\dox{\neg A}$。因此，我们必须先把事实 $D$ 加入行动前的模型 $M$，经由 $U$ 概括其影响，然后在修改 $\dox{\neg A}$ 发生之后重新求 $D$ 的值。于是，尽管行动的因果效应可以在句法上表述为反事实语句，这种“让已知事实的影响绕道 $U$”的需要，却使反事实成为与行动不同的一类语句（见 \S 1.4）。

我们还应强调，自然语言中多数反事实话语都（往往是隐含地）预设了关于受前件影响之事实的知识。例如，当我们说“若不是因为 $A$，$B$ 会有所不同”时，我们便隐含地知道 $B$ 的实际取值是什么，并且知道 $B$ 对 $A$ 是敏感的。正是这类关系赋予反事实以独特的品格——有别于行动语句——而且正如 \S 1.4 所见，也正是这类语句，其求值需要更详细的设定：必须具备关于函数机制 $f_i(\pa_i,u_i)$ 的某种知识。

\subsection{反事实的求值：概率分析}

为演示反事实的概率求值（式 \textup{(7.3)}—\textup{(7.5)}），我们对行刑队的故事稍作修改，假定：
\begin{enumerate}[label=\arabic*.,leftmargin=1.8em]
\item 法庭下达处决命令的概率为 $P(U)=p$；
\item 枪手 $A$ 因紧张而扣动扳机的概率为 $q$；且
\item 枪手 $A$ 的紧张与 $U$ 独立。
\end{enumerate}
在这些假定下，我们希望计算 $P(\neg D_{\neg A}\mid D)$——即在已知囚犯事实上已死的条件下，若 $A$ 没有开枪、囚犯仍然活着的概率。

直观上，我们可以这样得出答案：$\neg D_{\neg A}$ 为真当且仅当法庭没有下达命令。于是任务归结为计算 $P(\neg U\mid D)$，其值为 $q(1-p)/[1-(1-q)(1-p)]$。然而，我们的目的是演示一种基于式 \textup{(7.4)} 的、一般而形式的推导此类概率的方法，尽量少用直觉。

与新故事相联系的概率因果模型（定义 7.1.6）含两个背景变量 $U$ 与 $W$，其中 $W$ 表示枪手 $A$ 的紧张。该模型如下。

\begin{center}
\begin{tabular}{r@{\hspace{1.0em}}l@{\hspace{4.5em}}l}
\multicolumn{3}{c}{\textbf{模型 $\langle M,P(u,w)\rangle$}}\\[4pt]
 & $(U,W)\sim P(u,w)$ & \\[2pt]
$C=U$ & $(C)$ & \\
$A=C\vee W$ & $(A)$ & \\
$B=C$ & $(B)$ & \\
$D=A\vee B$ & $(D)$ &
\end{tabular}
\end{center}

在此模型中，背景变量的分布为
\begin{equation}
P(u,w)=
\begin{cases}
pq, & \text{若 } u=1,\ w=1,\\
p(1-q), & \text{若 } u=1,\ w=0,\\
(1-p)q, & \text{若 } u=0,\ w=1,\\
(1-p)(1-q), & \text{若 } u=0,\ w=0.
\end{cases} \tag{7.7}
\end{equation}

依定理 7.1.7，我们的第一步（溯因）是计算后验概率 $P(u,w\mid D)$，以计入“发现囚犯已死”这一事实。容易算得：
\begin{equation}
P(u,w\mid D)=
\begin{cases}
\dfrac{P(u,w)}{1-(1-p)(1-q)}, & \text{若 } u=1\ \text{或}\ w=1,\\[8pt]
0, & \text{若 } u=0\ \text{且}\ w=0.
\end{cases} \tag{7.8}
\end{equation}
第二步（行动）是构造子模型 $M_{\neg A}$，同时保留式 \textup{(7.8)} 的后验概率。

\begin{center}
\begin{tabular}{r@{\hspace{1.0em}}l@{\hspace{4.5em}}l}
\multicolumn{3}{c}{\textbf{模型 $\langle M_{\neg A},P(u,w\mid D)\rangle$}}\\[4pt]
 & $(U,W)\sim P(u,w\mid D)$ & \\[2pt]
$C=U$ & $(C)$ & \\
$\neg A$ & $(A)$ & \\
$B=C$ & $(B)$ & \\
$D=A\vee B$ & $(D)$ &
\end{tabular}
\end{center}

最后一步（预测）是在这个概率模型中计算 $P(\neg D)$。注意 $\neg D\iff\neg U$，于是结果（与预期一致）为
\[
P(\neg D_{\neg A}\mid D)=P(\neg U\mid D)=\frac{q(1-p)}{1-(1-q)(1-p)}.
\]

\subsection{孪生网络方法}

刚才所述程序的一个主要实践困难，在于需要计算、存储并使用后验分布 $P(u\mid e)$，其中 $u$ 表示模型中一切背景变量构成的集合。如前一例所示，即便我们出发时有一个背景变量彼此独立的马尔可夫模型，以 $e$ 为条件通常也会破坏这种独立性，从而不得不携带 $U$ 在给定 $e$ 下的联合分布的完整描述。若像我们在式 \textup{(7.8)} 中那样以表格形式编码，这种描述可能大到无法承受。

Balke and Pearl (1994b) 描述了一种克服此困难的图形方法：它使用两个网络，一个表示实际世界，一个表示假设世界。图 7.2 对所分析的行刑队故事展示了这一构造。

两个网络在结构上完全相同，唯一的差别是进入 $A^*$ 的箭头被删去，以映照从 $M_{\neg A}$ 中删去的方程。像暹罗连体婴一样，两个网络共享背景变量（在我们的例子中是 $U$ 与 $W$），因为它们在修改之下保持不变。内生变量则各复制一份并分别标号，因为它们在假设世界与实际世界中可能取不同的值。于是，计算模型 $\langle M_{\neg A},P(u,w\mid D)\rangle$ 中 $P(\neg D)$ 的任务，就归结为在所示的孪生网络中令 $A^*$ 为假、计算 $P(\neg D^*\mid D)$。

\begin{figure}[htbp]
\centering
\begin{tikzpicture}[causal,scale=1.0,every node/.style={transform shape}]
  % ---- 实际世界 ----
  \node[vtx] (u1) at (0,3.0) {$U$};
  \node[vtx] (w1) at (0.78,1.95) {$W$};
  \node[vtx] (c1) at (-0.25,1.75) {$C$};
  \node[vtx] (a1) at (-1.25,0.62) {$A$};
  \node[vtx] (b1) at (0.75,0.62) {$B$};
  \node[vtx] (d1) at (-0.25,-0.55) {$D$};
  % ---- 假设世界 ----
  \node[vtx] (c2) at (3.35,1.75) {$C^*$};
  \node[vtx] (a2) at (2.35,0.62) {$A^*$};
  \node[vtx] (b2) at (4.35,0.62) {$B^*$};
  \node[vtx] (d2) at (3.35,-0.55) {$D^*$};
  % ---- 边 ----
  \draw[ed] (u1)--(c1);
  \draw[ed] (u1)--(c2);
  \draw[ed] (w1)--(a1);
  \draw[ed] (c1)--(a1);
  \draw[ed] (c1)--(b1);
  \draw[ed] (a1)--(d1);
  \draw[ed] (b1)--(d1);
  \draw[ed] (c2)--(b2);
  \draw[ed] (a2)--(d2);
  \draw[ed] (b2)--(d2);
\end{tikzpicture}
\caption{行刑队的孪生网络表示。}
\label{fig:7.2}
\end{figure}

一般地，若我们希望计算反事实概率 $P(Y_x=y\mid z)$，其中 $X,Y,Z$ 是任意变量集（不必互不相交），则定理 7.1.7 指示我们在子模型 $\langle M_x,P(u\mid z)\rangle$ 中计算 $P(y)$；这归结为在一个增广贝叶斯网络中计算普通的条件概率 $P(y^*\mid z)$。这样的计算可以用标准的证据传播技术完成。把这一计算委托给贝叶斯网络中的推断，其好处在于：分布 $P(u\mid z)$ 无须显式写出，条件独立性可以被利用，并且可以采用局部计算方法（如 \S 1.2.4 所总结的那些）。

孪生网络表示还为检验反事实量之间的独立性提供了一个有用的办法。为作说明，设我们有一个链状因果图 $X\to Z\to Y$，并希望检验给定 $Z$ 时 $Y_x$ 是否与 $X$ 独立（即 $Y_x\indep X\mid Z$）。与这条链相联系的孪生网络如图 7.3 所示。

\begin{figure}[htbp]
\centering
\begin{tikzpicture}[causal,scale=1.0,every node/.style={transform shape}]
  \node[vtx] (ux) at (1.05,3.0) {$U_X$};
  \node[vtx] (x1) at (-0.55,2.45) {$X$};
  \node[vtx] (x2) at (2.05,2.45) {$X^*$};
  \node[vtx] (uz) at (1.05,1.45) {$U_Z$};
  \node[vtx] (z1) at (-0.55,0.95) {$Z$};
  \node[vtx] (z2) at (2.05,0.95) {$Z^*$};
  \node[vtx] (uy) at (1.05,0.05) {$U_Y$};
  \node[vtx] (y1) at (-0.55,-0.55) {$Y$};
  \node[vtx] (y2) at (2.05,-0.55) {$Y^*$};
  \draw[ed] (ux)--(x1);
  \draw[ed] (x1)--(z1);
  \draw[ed] (x2)--(z2);
  \draw[ed] (z1)--(y1);
  \draw[ed] (z2)--(y2);
  \draw[ed] (uz)--(z1);
  \draw[ed] (uz)--(z2);
  \draw[ed] (uy)--(y1);
  \draw[ed] (uy)--(y2);
\end{tikzpicture}
\caption{模型 $X\to Z\to Y$ 中反事实 $Y_x$ 的孪生网络表示。}
\label{fig:7.3}
\end{figure}

要检验 $Y_x\indep X\mid Z$ 是否在原模型中成立，我们检验在孪生网络中 $Z$ 是否把 $X$ 与 $Y^*$ d-分离。容易看出（依定义 1.2.3），以 $Z$ 为条件会使 $X$ 与 $Y^*$ 之间的路径经由 $Z$ 处的碰撞点而 d-连通，因此 $Y_x\indep X\mid Z$ 在该模型中并不成立。这一结论很难从链模型本身、或从该模型的方程中看出来。同理可见，每当我们以 $Y$ 或以 $\{Y,Z\}$ 为条件时，便在 $Y^*$ 与 $X$ 之间建立了联结；因此给定这些变量时，$Y_x$ 与 $X$ 并不独立。然而，若我们既不以 $Y$ 也不以 $Z$ 为条件，这一联结便被打断，此时 $Y_x\indep X$。

孪生网络揭示了形如 $Z_{\pa_Z}$ 的反事实的一种有趣解释，其中 $Z$ 是任一变量，$\PA_Z$ 表示 $Z$ 的父节点集。考虑如下问题：在图 7.3 的模型中，$Z_x$ 是否与某个给定变量集独立？此问题的答案取决于 $Z^*$ 是否与该变量集 d-分离。然而，任何与 $Z^*$ d-分离的变量，也必与 $U_Z$ d-分离；因此，表示 $U_Z$ 的节点可以充当反事实变量 $Z_x$ 的单向代理。这并非巧合：考虑到 $Z$ 由方程 $z=f_Z(x,u_Z)$ 支配，按定义，$Z_x$ 的概率等于在 $X$ 被固定于 $x$ 的条件下 $Z$ 的概率。在这种条件下，$Z$ 只有当 $U_Z$ 变化时才可能变化。因此，若 $U_Z$ 满足某种独立关系，则 $Z_x$（更一般地，$Z_{\pa_Z}$）也必满足该关系。于是，对形如 $Z_{\pa_Z}$ 的任何反事实变量，我们得到了一个以所谓“误差项”$U_Z$ 表述的简单图形表示。利用这一表示，我们可以容易地从图 7.3 验证 $(U_Y\indep X\mid\{Y^*,Z^*\})_G$ 与 $(U_Y\indep U_Z\mid\{Y,Z\})_G$ 在孪生网络中均成立，从而
\[
Y_z\indep X\mid\{Y_x,Z_x\}\qquad\text{与}\qquad Y_z\indep Z_x\mid\{Y,Z\}
\]
必在该模型中成立。关于孪生网络的更多考量——包括向多网络（表示不同前件下的反事实）的推广——见 Shpitser and Pearl (2007)，另见 \S 11.3.2 与 \S 11.7.3。

\section{结构模型的应用与解释}

\subsection{线性计量经济模型中的政策分析：一个例子}

在 \S 1.4 中，我们用需求与价格均衡这一经典经济问题（见图 7.4）例示了结构方程建模的性质。本章我们用它来回答与政策有关的问题。

回顾一下，这个例子由两个方程构成
\begin{equation}
q=b_1p+d_1i+u_1, \tag{7.9}
\end{equation}
\begin{equation}
p=b_2q+d_2w+u_2, \tag{7.10}
\end{equation}
其中 $q$ 是家庭对产品 $A$ 的需求量，$p$ 是产品 $A$ 的单位价格，$i$ 是家庭收入，$w$ 是生产产品 $A$ 的工资率，$u_1$ 与 $u_2$ 表示误差项——分别影响数量与价格的遗漏因素（Goldberger 1992）。

若定义 $V=\{Q,P\}$、$U=\{U_1,U_2,I,W\}$，并假定每个方程都表示定义 7.1.3 意义上的自主过程，则这个方程组构成一个因果模型（定义 7.1.1）。通常假定 $I$ 与 $W$ 可观测，而 $U_1$ 与 $U_2$ 不可观测且与 $I$、$W$ 独立。由于误差项 $U_1$ 与 $U_2$ 不可观测，模型的完整设定必须包含这些误差的分布，通常取为高斯分布，其协方差矩阵为 $\Sigma_{ij}=\mathrm{cov}(u_i,u_j)$。经济学中早已知道（可上溯至 Wright 1928）：线性性、正态性以及 $\{I,W\}$ 与 $\{U_1,U_2\}$ 的独立性这三项假定，足以一致地估计全部模型参数，包括协方差矩阵 $\Sigma_{ij}$。然而本书的关注点不是参数的估计，而是参数在政策预测中的运用。据此，我们将演示如何求值下列三个查询。

\begin{enumerate}[label=\arabic*.,leftmargin=1.9em,itemsep=0.3em]
\item 若价格被\emph{控制}在 $P=p_0$，需求量 $Q$ 的期望值是多少？
\item 若价格被\emph{报告}为 $P=p_0$，需求量 $Q$ 的期望值是多少？
\item 已知当前价格为 $P=p_0$，若我们\emph{把}价格\emph{控制}在 $P=p_1$，需求量 $Q$ 的期望值会是多少？
\end{enumerate}

读者应当认出，这三个查询分别代表行动、预测与反事实——即我们的三层级体系。第二个查询代表预测，在文献中是标准的，可以直接从协方差矩阵回答，无须涉及因果性、结构或不变性。第一个与第三个查询则立足于方程的结构性质，一如所料，标准的结构方程文献并不处理它们。\footnote{我曾把这一例子讲给全美一百多位计量经济学专业的学生与教师。听众回答第二个问题毫无困难，只有一人能解出第一个问题，无人能答出第三个问题。第五章（\S 5.1）给出了一种解释，\S 11.5.4 则依据 Heckman and Vytlacil (2007) 给出了更新的评估。}

为回答第一个查询，我们把式 \textup{(7.10)} 替换为 $p=p_0$，于是剩下
\begin{equation}
q=b_1p+d_1i+u_1, \tag{7.11}
\end{equation}
\begin{equation}
p=p_0, \tag{7.12}
\end{equation}
而 $U_1$ 与 $I$ 的统计特性保持不变。受控需求于是为 $q=b_1p_0+d_1i+u_1$，其期望值（以 $I=i$ 为条件）由
\begin{equation}
E[Q\mid \dox{P=p_0},i]=b_1p_0+d_1i+E(U_1\mid i) \tag{7.13}
\end{equation}
给出。由于 $U_1$ 与 $I$ 独立，末项等于
\[
E(U_1\mid i)=E(U_1)=E(Q)-b_1E(P)-d_1E(I),
\]
代回式 \textup{(7.13)} 便得
\[
E[Q\mid \dox{P=p_0},i]=E(Q)+b_1\bigl(p_0-E(P)\bigr)+d_1\bigl(i-E(I)\bigr).
\]

第二个查询的答案，由对式 \textup{(7.9)} 以当前观测 $\{P=p_0,\ I=i,\ W=w\}$ 为条件并取期望得到，
\begin{equation}
E(Q\mid p_0,i,w)=b_1p_0+d_1i+E(U_1\mid p_0,i,w). \tag{7.14}
\end{equation}
一旦给定 $\Sigma_{ij}$，$E[U_1\mid p_0,i,w]$ 的计算便是标准程序（Whittaker 1990, p.~163）。注意，尽管曾假定 $U_1$ 与 $I$、$W$ 独立，一旦观测到 $P=p_0$，这种独立性便不再成立。还要注意，式 \textup{(7.9)} 与 \textup{(7.10)} 都参与了求解，因而即便 $b_1=0$，观测值 $p_0$ 也会（通过 $E(U_1\mid p_0,i,w)$）影响期望需求 $Q$——这一点在查询 1 中并不成立。

第三个查询要求反事实量 $Q_{p=p_1}$ 的期望，以当前观测 $\{P=p_0,\ I=i,\ W=w\}$ 为条件（见 \S 11.7.1）。依定义 7.1.5，$Q_{p=p_1}$ 由子模型
\begin{equation}
q=b_1p+d_1i+u_1, \tag{7.15}
\end{equation}
\begin{equation}
p=p_1 \tag{7.16}
\end{equation}
支配；$u_1$ 的密度应以观测 $\{P=p_0,\ I=i,\ W=w\}$ 为条件。于是得到
\begin{equation}
E(Q_{p=p_1}\mid p_0,i,w)=b_1p_1+d_1i+E(U_1\mid p_0,i,w). \tag{7.17}
\end{equation}
期望值 $E(U_1\mid p_0,i,w)$ 与第二个查询的解相同；后者只在 $b_1p_1$ 这一项上有所不同。线性高斯模型中反事实查询求值的一般矩阵方法，见 Balke and Pearl (1995a)。

\begin{figure}[htbp]
\centering
\begin{tikzpicture}[causal,scale=1.0,every node/.style={transform shape}]
  \node[lat,minimum size=9mm] (U1) at (-3.2,1.9) {$U_1$};
  \node[vtx,minimum size=9mm] (I)  at (-1.55,2.15) {$I$};
  \node[vtx,minimum size=9mm] (W)  at (1.55,2.15) {$W$};
  \node[lat,minimum size=9mm] (U2) at (3.2,1.9) {$U_2$};
  \node[vtx,minimum size=9mm] (Q)  at (-1.85,0.3) {$Q$};
  \node[vtx,minimum size=9mm] (P)  at (1.85,0.3) {$P$};
  \draw[ed] (U1)--(Q);
  \draw[ed] (I)--(Q) node[midway,right,font=\small] {$d_1$};
  \draw[ed] (W)--(P) node[midway,right,font=\small] {$d_2$};
  \draw[ed] (U2)--(P);
  \draw[ed] (Q.north east) -- (P.north west) node[midway,above,font=\small] {$b_2$};
  \draw[ed] (P.south west) -- (Q.south east) node[midway,below,font=\small] {$b_1$};
\end{tikzpicture}
\caption{例示价格（$P$）与需求（$Q$）之关系的因果图。}
\label{fig:7.4}
\end{figure}

至此值得强调：反事实期望的计算问题并非学院式习题；事实上，它代表了几乎所有决策情境中的典型情形。每当我们着手预测政策的效应时，就有两项考量在起作用。第一，政策变量（如经济学中的价格与利率，过程控制中的压力与温度）很少是外生的。当我们观测一个运行中的系统时，政策变量是内生的；而到了规划阶段、当我们考虑各种行动与改变时，它们就变成外生的。第二，政策很少被抽象地评估；它们往往是被某些需要补救性纠正的偶发事件带到焦点上的。例如在故障排查中，我们观测到不受欢迎的效应 $e$，它受其他条件 $X=x$ 的影响；我们希望预测：采取某个使 $X$ 发生变化的行动，能否补救这一局面。$e$ 所提供的信息极有价值，在预测任何行动的效应之前，必须先（用溯因）处理这些信息。正是这一步溯因，赋予关于行动的实践性查询以反事实的品格，正如我们在第三个查询 \textup{(7.17)} 的求值中所见。

当前价格 $p_0$ 反映了决策之时盛行的经济条件（如 $Q$），而这些条件被推定为可以由所考虑的政策加以改变。因此，价格 $P$ 表示一个内生的决策变量（如图 7.4 所示），它在斟酌阶段按子模型 $M_{p=p_1}$ 的规定变为外生。查询 3 的假设语气，翻译成一个实践中的政策分析问题便是：“已知当前价格为 $P=p_0$，若我们今日把价格改为 $P=p_1$，求需求（$Q$）的期望值。”在实践决策情境中使用假设性措辞的理由，将在下一节以及 \S 11.7.2 中讨论。

\subsection{反事实的经验内容}

“反事实”一词是个误称，因为它暗示一种与事实相反的陈述，或者至少是一种逃逸于经验验证之外的陈述。反事实两者都不是；它们对科学思维至为根本，并且承载着与任何科学定律同样清晰的经验信息。

试看欧姆定律 $V=IR$。这一定律的经验内容可以用两种彼此替代的形式编码。
\begin{enumerate}[label=\arabic*.,leftmargin=1.9em,itemsep=0.35em]
\item \textbf{预测形式}：若在时刻 $t_0$ 我们测得电流 $I_0$ 与电压 $V_0$，则在其他条件不变之下，对任何未来时刻 $t\ge t_0$，若电流为 $I(t)$，则电压将为 $V(t)=\dfrac{V_0}{I_0}I(t)$。
\item \textbf{反事实形式}：若在时刻 $t_0$ 我们测得电流 $I_0$ 与电压 $V_0$，则假如时刻 $t_0$ 的电流是 $I'$ 而非 $I_0$，电压本会是 $V'=\dfrac{V_0}{I_0}I'$。
\end{enumerate}
表面上看来，预测形式作出了有意义且可检验的经验断言，而反事实形式只是在猜测未曾发生（也不可能发生）的事件，因为不可能在同一时刻把两个不同的电流通入同一个电阻。然而，若我们把反事实形式理解为不多不少正是预测形式的一种会话式缩写，前者的经验内容便清晰地透射出来。二者都使我们能从仅仅一次测量 $(I_0,V_0)$ 作出无穷多次预测；二者也都从一条科学定律获得其效力——该定律把某个不随时间变化的性质（比值 $V/I$）赋予任何导电物体。

但若反事实语句只不过是以迂回方式陈述一组预测，我们何不直接使用预测形式，而偏要诉诸如此繁复的表达方式？一个显然的答案是：我们常常用反事实来传达的不是预测本身，而是这些预测的逻辑后果。例如，说“假如 $A$ 当时没有开枪，囚犯本会还活着”，其意图可能只是传达“$B$ 没有开枪”这一事实信息。此时反事实的语气，只是为所传达的事实补上一个基于普遍定律的逻辑理由。一个不那么显然的答案，则在于预测性断言所附带的 ceteris paribus（其他条件不变）限定，而这个限定并非完全没有含混之处。当我们改变电阻中的电流时，究竟什么应保持不变——温度？实验设备？一天中的时刻？当然不会是电压表的读数！

当我们郑重地作出预测性断言时，这类事项必须仔细规定。而当我们使用反事实表达时，其中许多规定都是隐含的（因而多余），尤其是在我们对背后的因果模型已有共识的时候。例如，我们不必规定这些预测应在何种温度与压强下成立；这些已由“假如时刻 $t_0$ 的电流是 $I'$ 而非 $I_0$”这句话所蕴含。换言之，我们所指的恰恰是时刻 $t_0$ 时我们实验室中盛行的那些条件。这句话还蕴含：我们并不是真要让任何人把电压表的读数保持恒定；变量应循其自然进程演化，我们所设想的唯一改变，是那个（按我们的因果模型）当前决定着电流的机制。

总结起来，反事实语句完全可以被解释为：在一组明确规定的条件下——即该语句事实部分所盛行的条件——传达一组预测。为使这些预测有效，有两个成分必须保持不变：定律（或机制）与边界条件。用结构模型的语言说，定律对应于方程 $\{f_i\}$，边界条件对应于背景变量 $U$ 的状态。因此，反事实语句之预测式解释得以成立的一个前提，就是假定：当我们的预测性断言被应用或被检验时，$U$ 不会改变。

这一点最好用一个打赌的例子来说明。我们必须就一次公平掷硬币的结果押正面或反面；猜对了赢一美元，猜错了输一美元。假设我们押正面并且赢了一美元，却并未瞥一眼硬币的结果。考虑反事实“假如我当时押的是另一面，我就会输掉一美元”。这句话的预测式解释，翻译过来便是那个难以置信的断言：“假如我下一次押反面，我就会输掉一美元。”为使该断言有效，必须假定两个不变量：赔付规则与硬币的结果。前者在赌博情境中是合理的假定，后者却只在罕见的情形下才成立。正因如此，“假如我当时押的是另一面，我就会输掉一美元”这句话的预测效用相当低，有些人甚至视之为事后诸葛的废话。正是 $U$ 与 $f(x,u)$ 在时间上的持存，赋予反事实表达以预测力；一旦缺乏这种持存，反事实便失去其显而易见的预测效用。

然而，反事实中还有一层效用，并不直接转化为预测上的收益，因而或许可以解释反事实为何在人类话语中无处不在。我所想到的是解释价值。假设在那个打赌的故事里，每一次下注都重新掷币。“假如我当时押的是另一面，我就会输掉一美元”这句话就毫无价值了吗？我相信它有价值：它告诉我们，我们面对的并不是一个任性的庄家，而是一个至少会看一眼赌注、把它与某种标准作比较、并依一贯规则判定输赢的庄家。这一信息对作为赌客的我们也许用处不大，但对那些定期前来校准赌博机、以确保州政府抽成的检查员来说却可能有用。更重要的是，如果我们胆敢稍微作弊——比如说操纵硬币的轨迹，或安装一个微型发射器告诉我们硬币落在了哪一面——这一信息对我们赌客也可能有用。要让这类作弊奏效，我们就应当知道赔付规则 $y=f(x,u)$；而“假如我当时押的是另一面，我就会输掉一美元”这句话，恰好揭示了该规则的重要方面。

通过假设赌客作弊、规则被违反这类不大可能的情境来论证反事实的价值，是否过于牵强？我认为，这类不大可能的操作恰恰是衡量语句解释价值的常规做法。任何因果解释的本性就在于：它的效用不是在标准情境中得到证明，而是在需要对新标准作创新性操作的新颖情境中得到证明。理解电视机如何工作的效用，并不来自正确地转动旋钮，而来自电视坏掉时能够把它修好。回想一下，每个因果模型所宣告的不是一个子模型，而是一大批子模型，每一个都因违反某些定律而被创造出来。于是，因果模型中机制的自主性，便意味着一种公开的邀请：去移除或替换那些机制；而语句的解释价值以“它们能多好地预测此类替换的后果”来衡量，也就再自然不过了。

\subsubsection*{带内在非确定性的反事实}

回顾此前的讨论，我们看到反事实可在两种条件下获得预测价值：(1) 当产生不确定性的未观测变量（$U$）保持不变（直到我们下一次预测或行动）；或 (2) 当产生不确定性的变量具有在未来某个时刻（在我们下一次预测或行动之前）被观测到的可能。两种情形下，我们都还需确保产生结果的机制 $f(x,u)$ 保持不变。

这些结论引出了一些有趣的问题，涉及反事实在微观现象中的使用：对于我们在量子论中遇到的那类不确定性，上述条件一个也不成立。海森堡的骰子每秒钟被重新掷出数十亿次，而我们对 $U$ 的测量永远不够精细，无法从响应方程 $y=f(x,u)$ 中消除一切不确定性。因此，当我们在分析中纳入量子层面的过程时，便面对一个两难：要么放弃一切关于反事实的谈论（这是包括 Dawid 2000 在内的一些研究者所建议的策略），要么继续使用反事实、但把其使用限于它们具有经验意义的情境。这等于说，在我们的分析中只保留那些满足上一段条件 (1) 与 (2) 的 $U$。我们不再假设能完全消除一切不确定性的 $U$，而只承认那些或者 (1) 持存、或者 (2) 可能被观测到的 $U$。

自然，粗化背景变量的粒度是有代价的：机制方程 $v_i=f_i(\pa_i,u_i)$ 失去其确定性品格，因而应被改造为随机式的。我们不应再由一组确定性方程 $\{f_i\}$ 构造因果模型，而应考虑由随机函数 $\{f_i^*\}$ 构成的模型，其中每个 $f_i^*$ 是从 $V\cup U$ 到 $V_i$ 各状态上某个内在概率分布 $P^*(y_i)$ 的映射。这一选择导向因果贝叶斯网络（\S 1.3）：其中条件概率 $P^*(y_i\mid \pa_i,u_i)$ 表示内在非确定性（intrinsic nondeterminism，有时称为“客观概率”，objective chance；Skyrms 1980），而根节点表示或者持存、或者可能被观测到的背景变量 $U$。在这一表示中，反事实概率 $P(Y_x=y\mid e)$ 仍可用定理 7.1.7 的三步（溯因、行动、预测）求值。在溯因阶段，我们以可用的证据 $e$ 为条件，更新根节点的先验概率 $P(u)$，得到 $P(u\mid e)$。在行动阶段，我们删去进入集合 $X$ 中变量的箭头，并把它们的取值实例化为 $X=x$。最后，在预测阶段，我们计算由更新后的受操纵网络所产生的 $Y=y$ 的概率。

这一求值当然也可以在普通的因果贝叶斯网络中实现（即不限于表示内在非确定性的网络），但那样算得的结果并不表示反事实 $Y_x=y$ 的概率。那样的求值等于假定个体是同质的，每个个体都具备总体的随机性质——即 $P(y_i\mid \pa_i,u)=P(y_i\mid \pa_i)$。这种假定在量子层面的现象中或许是恰当的（那里个体代表特定的实验条件），但在宏观现象中则不恰当，因为个体之间可能差异显著。在第一章的例子中（\S 1.4.4，图 1.6），这种随机化归因等于假定没有任何个体受药物影响（如模型 1 所规定），而忽略了某些个体事实上可能比另一些对药物更敏感的可能性（如模型 2）。

\subsection{因果解释、语句及其诠释}

“解释增进理解，而理解得越多的人推理与学习越有效”，这是老生常谈。人们也普遍接受：解释这一概念无法与因果关系相分离；例如，症状可以解释我们对某种疾病的信念，却不解释疾病本身。然而，原因与解释之间的确切关系仍是众说纷纭的话题（Cartwright 1989; Woodward 1997）。拥有一个在确定性与概率性两种设定下都成立的形式化因果性与反事实理论，会为“什么构成一个恰当的解释”这一问题带来新的启示，并为机器自动生成解释开辟新的可能（Halpern and Pearl 2005a,b）。

生成解释的一个自然起点，是使用因果贝叶斯网络（\S 1.3）：其中待解释的事件（被解释项，explanandum）由网络中被实例化节点的某种组合 $e$ 构成，而任务是找出 $e$ 的祖先（即原因）之某个子集的一个实例化 $c$，使某种“解释力”的度量最大化——所谓解释力，即 $c$ 解释 $e$ 的程度。然而，这一度量的恰当选择尚无定论。许多哲学家与统计学家主张以似然比
\[
L=\frac{P(e\mid c)}{P(e\mid c')}
\]
作为“$c$ 比 $c'$ 更好地解释 $e$”之程度的恰当度量。在 Pearl (1988b, chap.~5) 与 Peng and Reggia (1986) 中，最佳解释则通过最大化后验概率 $P(c\mid e)$ 求得。这两种度量各有缺陷，并已受到若干研究者的批评，包括 Pearl (1988b)、Shimony (1991, 1993)、Suermondt and Cooper (1993) 与 Chajewska and Halpern (1997)。为补救这些缺陷，人们提出过对概率参数 $[P(e\mid c),P(e\mid c'),P(c),P(c')]$ 作更精细的组合，但似乎没有一种能很好地把握人们赋予“解释”一词的含义。

概率度量的毛病在于，它们无法刻画 $c$ 与 $e$ 之间因果联系的强度；任何命题 $h$，只要稍作想象，都可以被认为对 $e$ 有某种影响，无论多么微弱。这便使 $h$ 有资格成为因果网络中 $e$ 的祖先，并凭借 $h$ 与 $e$ 之间强烈的虚假关联而与真正的解释相竞争并胜出。

要摆脱这一困难，我们必须越出概率度量，转而专注于因果参数，例如因果效应 $P(y\mid\dox{x})$ 与反事实概率 $P(Y_{x'}=y'\mid x,y)$，以之作为定义解释力的基础。这里 $x$ 与 $x'$ 取遍各个备择解释，$Y$ 是被观测到取值 $y$ 的那组响应变量。表达式 $P(Y_{x'}=y'\mid x,y)$ 读作：假如 $X$ 曾为 $x'$（而非实际取值 $x$），$Y$ 会取另一值 $y'$ 的概率。（注意 $P(y\mid\dox{x})\triangleq P(Y_x=y)$。）如今，用于计算因果效应与反事实概率的计算模型已发展起来，这使我们能把上述参数与标准的概率参数结合起来，从而综合出一个更忠实的解释力度量，用以指导恰当解释的选取与生成。

这些可能性触发了一个重要的基本问题：“解释”究竟是一个基于普遍原因（如“饮毒芹导致死亡”）的概念，还是基于单一原因（如“苏格拉底饮毒芹导致了他的死亡”）的概念？因果效应表达式 $P(y\mid\dox{x})$ 属于第一类，反事实表达式 $P(Y_{x'}=y'\mid x,y)$ 则属于第二类，因为以 $x$ 与 $y$ 为条件，把世界情景收窄到与手头最具体的信息 $X=x$、$Y=y$ 相容的那些。

把因果陈述分为普遍与单一两类，一直是哲学中大量研究的主题（例如见 Good 1961; Kvart 1986; Cartwright 1989; Eells 1991；另见 \S 7.5.4 与 \S 10.1.1 的讨论）。这类研究在认知科学与人工智能中却少有问津，部分因为它没有带来实际的推断程序，部分因为它立足于有问题的概率语义（关于概率因果性的讨论见 \S 7.5）。在机器生成解释的语境中，这一分类同时具有认知意义与计算意义。我们在第一章（\S 1.4）讨论过两类因果查询之间的鲜明分界线：一类可由 $\langle P(M),G(M)\rangle$（分别为与模型 $M$ 相联系的图与概率）回答，另一类则需要以函数设定的形式补充信息。普遍的因果陈述（如 $P(y\mid\dox{x})$）往往属于第一类（如第三章所示），而反事实表达式（如 $P(Y_{x'}=y'\mid x,y)$）属于第二类，因而要求更详细的设定与更高的计算资源。

把解释恰当地归入普遍类或单一类，取决于原因 $c$ 究竟是凭借其产生 $e$ 的普遍倾向（与 $c$ 的各备择项较弱的倾向相比）而获得相对于效应 $e$ 的解释力，还是凭借 $c$ 对于触发某条在当下具体情境中导致 $e$ 的事件链（该情境由 $e$ 以及也许其他事实与观测所刻画）乃是必要的。形式地说，差别取决于：在评估各个假设的解释力时，我们是否应把信念建立在确实发生的事件 $c$ 与 $e$ 之上。

对这些备择方案的形式分析见第九、十章，那里我们讨论因果关系的必要与充分方面，以及单一事件因果的概念。本节余下部分将关注解释性语句的诠释与生成，并以“必要”这一方面为规范。

下面这份清单主要取自 Galles and Pearl (1997)，它给出解释性话语中所用语句的例子，以及它们在 \S 7.1.1 所述可修改结构模型进路中的相应语义。

\begin{itemize}[leftmargin=1.6em,itemsep=0.35em]
\item “$X$ 是 $Y$ 的原因”，若存在 $X$ 的两个取值 $x$ 与 $x'$、以及 $U$ 的一个取值 $u$，使得
\[
Y_x(u)\ne Y_{x'}(u).
\]
\item “在语境 $Z=z$ 中，$X$ 是 $Y$ 的原因”，若存在 $X$ 的两个取值 $x$ 与 $x'$、以及 $U$ 的一个取值 $u$，使得
\[
Y_{xz}(u)\ne Y_{x'z}(u).
\]
\item “$X$ 是 $Y$ 的\emph{直接}原因”，若存在 $X$ 的两个取值 $x$ 与 $x'$、以及 $U$ 的一个取值 $u$，使得
\[
Y_{xr}(u)\ne Y_{x'r}(u),
\]
其中 $r$ 是 $V\backslash\{X,Y\}$ 的某个取值。
\item “$X$ 是 $Y$ 的\emph{间接}原因”，若 $X$ 是 $Y$ 的原因，且 $X$ 不是 $Y$ 的直接原因。
\item “事件 $X=x$ \emph{总是}导致 $Y=y$”，若：
\begin{enumerate}[label=(\roman*),leftmargin=2.0em,itemsep=0.1em]
\item 对一切 $u$ 有 $Y_x(u)=y$；且
\item 存在 $U$ 的一个取值 $u'$，使得对某个 $x'\ne x$ 有 $Y_{x'}(u')\ne y$。
\end{enumerate}
\item “事件 $X=x$ \emph{可能}导致了 $Y=y$”，若：
\begin{enumerate}[label=(\roman*),leftmargin=2.0em,itemsep=0.1em]
\item $X=x$ 与 $Y=y$ 为真；且
\item 存在 $U$ 的一个取值 $u$，使得 $X(u)=x$、$Y(u)=y$，且对某个 $x'\ne x$ 有 $Y_{x'}(u)\ne y$。
\end{enumerate}
\item “未观测事件 $X=x$ 是 $Y=y$ 的\emph{很可能的原因}（likely cause）”，若：
\begin{enumerate}[label=(\roman*),leftmargin=2.0em,itemsep=0.1em]
\item $Y=y$ 为真；且
\item 对一切 $x'\ne x$，$P(Y_x=y,\,Y_{x'}\ne y\mid Y=y)$ 都很高。
\end{enumerate}
\item “事件 $Y=y$ \emph{尽管} $X=x$ 仍然发生了”，若：
\begin{enumerate}[label=(\roman*),leftmargin=2.0em,itemsep=0.1em]
\item $X=x$ 与 $Y=y$ 为真；且
\item $P(Y_x=y)$ 很低。
\end{enumerate}
\end{itemize}

上述清单展示了可修改结构模型在形式化因果表达之细微差别方面的灵活性。更多的细微差别（涉及促成、阻止、维持、产生等概念）将在第九、十章分析。相关的表达还有：“事件 $A$ 解释了事件 $B$ 的发生”；“若 $C$ 成立，$A$ 便可解释 $B$”；“$B$ 是尽管 $A$ 而发生的，因为 $C$ 为真”。能够诠释并生成这类解释性语句，或者为语境挑选最恰当的表达，是人机对话研究中最引人入胜的挑战之一。

\subsection{从机制到行动再到因果}

\S 7.1.1 所述的结构模型语义，为认知科学与人工智能中的两个问题提示了解法：行动的表征，以及因果排序的作用。我们依次讨论这两个问题，因为后者建立在前者之上。

\subsubsection*{行动、机制与手术}

无论我们采取“行动是概率分布到概率分布的变换”这一概率范式，还是采取“行动是状态到状态的变换”这一确定性范式，这类变换在原则上都可能是无限复杂的。然而在实践中，人们相当快地就能互相教会世界上各种行动的通常结果，人们也能毫不费力地预测大多数行动的后果。何以如此？

结构模型对此的回答是：假定我们在日常推理中通常援引的行动，都可以表示为\emph{局部手术}（local surgery）。世界由数量巨大的自主而不变的联结或机制构成，每一个机制对应一个物理过程，约束着相对少数几个变量的行为。如果我们理解这些联结如何相互作用（通常它们只是共享变量），我们也就应当能够理解任一给定行动的效应会是什么：只需重新规定被该行动扰动的那少数几个机制；然后让修改后的机制组合相互作用，看系统在平衡态会演化到哪个状态。若设定是完备的（即 $M$ 与 $U$ 已给定），则演化为单一状态。若设定是概率性的（即 $P(u)$ 已给定），则浮现出一个新的概率分布；若设定是部分的（即某些 $f_i$ 未给定），则创造出一个新的、部分的理论。三种情形下，我们都应能够回答关于行动后事态的查询，只不过精确程度依次递减。

使这一方案可操作的关键成分，是行动的\emph{局部性}。孤立地看，局部性是个含混的概念，因为在一个空间中局部的东西，在另一个空间中未必局部。例如，一粒尘埃在频率（或傅里叶）表示中显得极为弥散；反过来，一个纯音则需要一段很长的时间才能被欣赏。结构语义强调：行动是在\emph{机制}的空间中局部，而不是在变量、语句或时间槽的空间中局部。例如，推倒一排多米诺骨牌中最左边的一块，在物理空间中似乎并不“局部”，然而它在机制域中却相当局部：只有一个机制被扰动，即通常使那块骨牌保持稳定直立的重力恢复力；其他一切机制都按设定保持不变，服从通常的物理方程。局部性使我们可以轻松地规定这一行动，而无须枚举其全部后果。听者若假定她与我们共享关于多米诺物理的理解，便能自行算出这一行动——或任何形如“把第 $i$ 块骨牌向右推倒”的行动——的后果。通过把领域表示为稳定机制的组合，我们实际上创造了一个“神谕”，它能回答关于巨大一整套行动及行动组合之效应的查询——而无须我们把这些效应逐一明确说出来。

\subsubsection*{定律与事实}

这套手术程序在结构方程模型的语境中表述出来时，听上去平淡无奇。然而，当人们试图在经典逻辑中实现此类方案时，却遇到了巨大的困难。要在机制空间中施行手术程序，我们需要一种语言，其中某些语句被赋予与其他语句不同的地位。描述机制的语句应当与描述生活中其他事实（如观测、假定与结论）的语句区别对待，因为前者被推定为稳定的，而后者是暂时的。的确，描述多米诺骨牌如何相互作用的方程保持不变，尽管骨牌自身的状态可以随境况而变。

承认这一区分的必要性，在逻辑进路处理行动与因果性的历程中曾是一次艰难的转变，也许因为经典逻辑的许多力量正来自其表征上的齐一性与句法上的不变性——其中没有哪条语句享有特殊地位。概率论者则远不那么抗拒接受定律与事实的区分，因为这一区分早在 1763 年就已被 Bayes 牧师编入概率语言：事实被表述为普通命题，因而可以获得概率值、可以被取条件；定律则被表述为条件概率语句（如 $P(\text{事故}\mid\text{疏忽驾驶})=$ 高），因而不能被赋予概率、也不能被取条件。正是由于这一传统，概率论者一向赋予条件句（如“鸟会飞”）非命题的品格，拒绝允许嵌套的条件句（Levi 1988），并坚持把人们对条件句的信心解释为条件概率判断（Adams 1975；另见 Lewis 1976）。值得注意的是，这些被某些哲学家视为局限的约束，恰恰是使概率论者不致混淆定律与事实的保障，使他们免于陷入某些曾困住逻辑进路的陷阱。\footnote{定律与事实的区分由 Poole (1985) 与 Geffner (1992) 提出，作为非单调推理的一条基本原理。在数据库理论中，定律由称为完整性约束（integrity constraint）的特殊语句表述（Reiter 1987）。随着人工智能中表述行动的需要日益明确，这一区分似乎正获得更广泛的支持（Sandewall 1994; Lin 1995）。}

\subsubsection*{机制与因果关系}

从目前的讨论看，似乎不诉诸任何因果概念，也能构造出计算行动后果的有效表示。在物理与工程的许多领域中的确如此。例如，若我们有一个由电阻与电压源构成的大型电路，并希望计算改变电路中某一个电阻的效应，则因果性概念几乎不进入计算。我们只需把电阻的修改值代入欧姆方程与基尔霍夫方程，然后求解这组（对称的）方程以得到所需变量。这一计算可以有效地完成，而无须对电流与电压之间的任何有向因果关系作出承诺。

要理解因果性的作用，我们应当注意：（与我们的电路例子不同）大多数机制在日常语言中并没有名字。我们说“加税”“逗他笑”“按下按钮”——一般地，$\doxp(q)$，其中 $q$ 是一个命题，而不是一个机制。在电路例子中说“增大这个电流”或“若这个电流更高……”会毫无意义，因为（最小地）增大该电流有许多方式，每一种都有不同的后果。显然，常识知识并不像电阻网络那样纠缠。再举一例：在 STRIPS 语言（Fikes and Nilsson 1971）中，行动不是由它所修改的机制之名来刻画的，而是由该行动的直接效果（ADD 表与 DELETE 表）来刻画的，而这些效果被表述为普通命题。的确，若我们的知识是按因果组织的，那么这种设定便是充分的，因为每个变量都由且仅由一个机制支配（见定义 7.1.1）。于是，我们应当能够自行判断：为实现所规定的效果，必须扰动的是哪一个机制；而这应当使我们能够预测情景的其余部分。

这种语言上的缩略定义出事件之间的一种新关系，我们通常称之为“因果关系”：若为实现 $A$ 所需的扰动蕴含 $B$ 的实现，则事件 $A$ 导致 $B$。\footnote{“所需”一词含有极小性之意，可译为：“……若每一个实现 $A$ 的极小扰动都蕴含 $B$。”这一蕴含关系的必要与充分方面在第九章（\S 9.2）中形式化。}这类因果缩略在表述领域知识时极为有效。关于哪些关系是稳定的、机制之间如何相互作用的复杂描述，很少以机制的语言被显式交流；它们往往以事件或变量之间因果关系的形式被交流。例如我们说：“若第 $i$ 块骨牌被向右推倒，它便导致第 $i+1$ 块骨牌也被向右推倒”；我们并不以每块骨牌维持其物理形状、响应重力牵引、服从牛顿力学的倾向来交流这类知识。

\subsection{西蒙的因果排序}

能够直接用“一个事件导致另一个事件”来谈论（而不是“一个行动改变了一个机制，而该改变又产生了效果”），在计算上非常有用；但与此同时，它要求我们领域中的机制组合满足某些条件，以容纳因果方向性。的确，\S 7.1.1 给出的因果模型形式定义假定：每个方程都被指定了一个独特的、位于其左端的特权变量，被视为“因变量”或“输出”。然而一般说来，一个机制可以被规定为函数约束
\[
G_k(x_1,\dots,x_l;\,u_1,\dots,u_m)=0,
\]
而不指明任何“因变量”。

Simon (1953) 设计了一个程序，用以判定一组这样的对称约束 $G$ 是否规定了为每个机制选取一个内生因变量的唯一方式（背景变量除外，因为它们由系统之外决定）。Simon 问道：我们何时能把变量 $(V_1,V_2,\dots,V_n)$ 排序，使得我们能解出每个 $V_i$ 而无须解出 $V_i$ 的任何后继？这样的排序若存在，便规定了我们赋予因果性的方向。这一准则初听之下颇显人为，因为解方程的顺序只是计算上的便利问题，而因果方向性却是物理实在的客观属性。（关于这一问题的讨论见 De Kleer and Brown 1986; Iwasaki and Simon 1986; Druzdzel and Simon 1993。）为给该准则辩护，让我们用行动与机制的语言重述 Simon 的问题。假定每个机制（即方程）都可以独立于其他机制而被修改，并令 $A_k$ 为能够修改方程 $G_k$（同时使其他方程不变）的那组行动。设想我们从 $A_k$ 中选出一个行动 $a_k$，并以这样的方式修改 $G_k$：整个方程组之解集 $(V_1(u),V_2(u),\dots,V_n(u))$ 不同于行动之前。若 $X$ 是被 $G_k$ 直接约束的那组变量，我们可以问：$X$ 中是否有某个成员，比方说 $X_k$，能够解释所有其他解的变化？若这一“起预测作用的”成员的身份，对一切 $a_k$ 与 $u$ 的选择都保持不变，我们便把 $X_k$ 指定为 $G_k$ 中的因变量。

形式地说，这一性质意味着：$a_k$ 的变化诱导出从 $X_k$ 的定义域到 $V\backslash X_k$ 的定义域的一个函数映射，即系统中（由 $a_k$ 产生的）一切变化都可归因于 $X_k$ 的变化。在这种情形下，把 $X_k$ 指定为机制 $G_k$ 的“代表”便是合理的；我们也就有理由把语句“行动 $a_k$ 导致事件 $Y=y$”替换为“事件 $X_k=x_k$ 导致 $Y=y$”（其中 $Y$ 是系统中任一变量）。$X_k$ 对 $a_k$ 之选择的这种不变性，正是把行动视为模态 $\dox{X_k=x_k}$ 的依据（定义 7.1.3）。它为我们以行动的即时后果来刻画该行动提供了许可，而不必顾及实际带来这些后果的工具；事实上，它定义的就是“局部行动”或“局部手术”这一概念。

可以证明（Nayak 1994）：$X_k$ 的唯一性可由一个只涉及方程组纯拓扑性质（即变量如何被归入各方程）的简单准则来判定。该准则是：应当能够在方程与变量之间建立一一对应，且该对应是唯一的。这可以通过求解方程与变量之间的“匹配问题”（matching problem；Serrano and Gossard 1987）来判定。若匹配唯一，则每个方程中因变量的选择便唯一，而该选择所诱导的方向性定义了一个有向无环图（DAG）。例如在图 7.1 中，箭头的方向性无须从外部指定；它可以由刻画该问题的对称约束（即逻辑命题）之集合
\begin{equation}
S=\{G_1(C,U),\ G_2(A,C),\ G_3(B,C),\ G_4(A,B,D)\} \tag{7.18}
\end{equation}
机械地确定。读者可以容易地验证：从每个方程中选出一个特权变量的方式是唯一的，因而图 7.1 所示箭头的因果方向性是不可避免的。

由此可见，按 Simon 的看法，因果方向性从两个假定中浮现：(1) 变量被划分为背景集（$U$）与内生集（$V$）；(2) 模型中机制的整体配置。据此，在某个机制中被指定为“因变量”的变量，当同一机制被嵌入另一个模型时，完全可能被标为“自变量”。的确，火车上坡时引擎导致车轮转动，下坡时二者的角色就变了。

当然，若我们无法确定背景变量，便可能引出若干种因果排序。例如在式 \textup{(7.18)} 中，若没有给定“$U$ 是背景变量”这一信息，则 $\{U,A,B,C\}$ 中任何一个都可被选作背景，而每一种选择都会在其余变量上诱导出不同的排序。（某些排序会与常识相冲突，例如“队长的信号影响法庭的决定”。）然而，$A\to D$ 的方向性在所有这些排序中都会保持。至于是否存在变量的一个划分 $\{U,V\}$，能使对称约束系统产生因果排序，这一问题同样可以用拓扑方法（在多项式时间内）解决（Dechter and Pearl 1991）。

当我们无法一次解出一个方程、而必须同时求解 $k$ 个方程构成的一块时，Simon 的排序准则便失效了。此时，由该块决定的全部 $k$ 个变量相互之间无序，尽管它们与其他块的关系仍可能有序。例如，这发生在图 7.4 的经济模型中：式 \textup{(7.9)} 与 \textup{(7.10)} 需要关于 $P$ 与 $Q$ 联立求解，因而方程与变量之间的对应不唯一；$Q$ 或 $P$ 都可以在两个方程中的任一个里被指定为“自变量”。的确，把式 \textup{(7.9)} 归类为“需求”方程（并相应地把式 \textup{(7.10)} 归为“价格”方程）所需的信息，并非来自变量被分派到各方程的方式，而是来自主题事务的考量。我们关于“家庭收入直接影响家庭需求（而非价格）”的理解，在这一归类中起了主要作用。

在那些我们倾向于断然宣称反馈回路中因果流沿顺时针方向的场合，这一断言通常立足于各种力在量级上的相对大小。例如，拧水龙头会降低水箱中的水位，但我们几乎无法对箱中之水做任何事来拧动水龙头。当此类信息可用时，因果方向性便再次通过诉诸假设性干预的概念来确定：追问对机制中某一变量的外部控制是否必然影响其他变量。这一考量于是构成识别非递归因果模型（定义 7.1.1）中因变量 $V_i$ 的操作语义。

刻画因果关系的这种不对称性，与物理方程的对称性毫不冲突。当我们说“$X$ 导致 $Y$，而 $Y$ 不导致 $X$”时，我们的意思是：改变一个以 $X$ 为通常因变量的机制，与改变一个以 $Y$ 为通常因变量的机制，二者对世界的影响不同。由于涉及两个不同的机制，这一陈述与我们在物理方程中所见的对称性完全协调。

Simon 的因果排序理论对休谟的因果归纳问题——即因果知识如何从经验中获得（见第二章）——有着深远的影响。能够从一组对称机制（连同对内生变量集的选择）推出因果方向性，这意味着：因果关系的获得与普通物理定律（如悬挂弹簧的胡克定律或牛顿的加速度定律）的获得（例如通过实验）并无二致。这并不意味着获得物理定律是件轻而易举、毫无方法论与哲学微妙之处的事情。它确实意味着：因果归纳这一哲学史上最棘手的问题之一，可以被化归为更为熟悉的科学归纳问题。
